\scanpage{55}
4.2. Définition des sites et des topos cristallins \inserted{de $X$ sur $S$}

Définition des sites et du topos stratifiant de $X$ sur $S$

C'est pareil si $X$ form$^{\mathrm{t}}$ lisse sur $S$.

\inserted{Ce sont des topos annelés}
\inserted{Attention, les sites \uline{n'ont pas en général d'objet final}!}
\deleted{Description d'un}
Ce sont des \deleted{\unsure{sites}} et topos \uline{annelés}.
Description des \deleted{un} faisceaux d'ensembles en termes de faisceaux habituels. Id. pour faisceaux d'anneaux, faisceaux de Modules sur faisceaux d'anneaux etc.

Faisceaux spéciaux: les hom de changement de base sont des isom.

\diagram{\[\begin{array}{l}
\text{Faisceaux de Modules spéciaux sur sites stratifiants}\\
\hfill\text{de }X/S\\
\rotatebox{90}{$\simeq$}\\
\text{Faisceaux de Modules à stratification sur }X/S.
\end{array}\]}

4.3. Énonçons de suite les résultats et conjectures suivants, qui seront \deleted{\unsure{établis}} détaillés plus tard:

\uline{Th} Si $S$ de car. 0 et $X$ lisse sur $S$, alors
\[
H^*_{\substack{\unsure{\mathrm{cris}/S}\\\text{ou strat}}}
(X/S,\underline{\mathcal{O}}_X)\simeq H^*_{\mathrm{DR}}(X/S)
\]
(isom can).
\deleted{\unsure{on a} $\unsure{f}=\mathbb{R}f_*(\mathcal{O}_X)$ \unsure{puis} $\uline{\unsure{H^i=0}}$, et
\[
(\mathbb{R}\unsure{f_{\mathrm{str}*}|f_*}(\mathcal{O}_X)^{\unsure{\mathrm{an}|\mathrm{ab}}}
\simeq\mathbb{R}f_*(\Omega^{\unsure{\bullet}}_{X/S})\quad ....]
\]}

\margin{54}
